\section{Related work}
\paragraph{World models and predictive state.}
World models learn representations that support forecasting and control
\citep{ha2018worldmodels,hafner2019planet,hafner2020dreamer}.
Predictive state representations likewise describe a system through predictions
of future observations \citep{littman2002predictive}.
SpikeWorld uses prediction in two phases: offline to train sensory and action
features, and online to obtain outcome-based correction signals. It stores those corrections outside the trained core.

\paragraph{Spiking temporal representations.}
Surrogate gradients enable optimization through discrete spike thresholds
\citep{neftci2019surrogate,zenke2018superspike}, and can support learning from
spike timing \citep{yu2026beyondrate}. Spiking Transformers combine spike-based
processing with attention \citep{zhou2023spikformer}. Event-based sensing
provides temporal contrast streams \citep{gallego2022event}; our experiments
also use encoded feature sequences from conventional data sources. We study how the shared core predicts these sequences and specify each
stream's input encoding and target.

\paragraph{Online adaptation and compact memory.}
Test-time training uses self-supervised deployment signals to update a model
\citep{sun2020ttt}. Fast weights and differentiable plasticity separate rapidly
changing information from slower learned parameters
\citep{ba2016fastweights,miconi2018plasticity}. SpikeWorld implements a related
separation through a fixed-size residual interface and a separate candidate
score state. The model's original semantic computation remains outside these
online writes.

\paragraph{Adaptive control.}
Online system identification and multiple-model adaptive control estimate
changes and adjust commands \citep{narendra1997multiple,yu2017uposi}.
Learned motor adaptation can infer context across changing robot dynamics
\citep{kumar2021rma}. SpikeWorld's candidate-score path belongs to this
identification-and-correction setting. Its residual matrices add a distinct
learned prediction correction, whose contribution is isolated experimentally.
The appendix compares corrections applied in actuator space and at the
prediction output.
